\title{QLoRA: Efficient Finetuning of Quantized LLMs}

\begin{abstract}
We introduce \textsc{QLoRA}, a finetuning method optimized for efficiency, enabling the adaptation of a 65B parameter model on a single 48GB GPU with no loss in full 16-bit finetuning performance. In \textsc{QLoRA}, gradients are computed through a 4-bit quantized, frozen pretrained language model, and propagated into Low Rank Adapters (LoRA). Our leading suite of models, termed \textbf{Guanaco}, surpasses all previously released open-source models on the Vicuna benchmark, achieving 99.3\% of ChatGPT’s performance while requiring only 24 hours of finetuning on one GPU. Key contributions of \textsc{QLoRA} for reducing memory consumption without compromising accuracy include: (a) the introduction of 4-bit NormalFloat (NF4), a novel datatype theoretically optimal for weights with normal distributions; (b) Double Quantization, which further lowers the average memory demands by quantizing the quantization coefficients themselves; and (c) Paged Optimizers that efficiently handle memory surges. Utilizing \textsc{QLoRA}, we finetune over 1,000 models and systematically analyze instruction-following and chatbot capabilities across 8 instruction datasets, a range of architectures (LLaMA, T5), and larger scales—such as 33B and 65B parameters—impractical for standard finetuning. Findings indicate that finetuning with QLoRA on compact, high-quality datasets delivers state-of-the-art outcomes, even when deploying models smaller than prior SoTA. We provide an extensive assessment of chatbot effectiveness using both human judgments and GPT-4-based evaluations, concluding that GPT-4 is an affordable and valid proxy for human assessment. Additionally, we observe that current chatbot benchmarks are unreliable indicators of true chatbot ability. Our targeted analysis highlights specific limitations of \textbf{Guanaco} relative to ChatGPT. We make available all models, code, and CUDA kernels necessary for 4-bit training.\footnote{\url{https://github.com/artidoro/qlora} and \url{https://github.com/TimDettmers/bitsandbytes}}
\end{abstract}

\section{Introduction}

Adjusting large language models (LLMs) via finetuning has proven to be an effective approach for enhancing their capabilities \citep{min2021metaicl, wei2021finetuned, ouyang2022training, wang2022super, wang2022self, liu2022few}, as well as for encouraging preferred behaviors or mitigating unwanted ones \citep{ouyang2022training,askell2021general,bai2022training}. Despite its advantages, finetuning massive models poses substantial computational challenges: for example, standard 16-bit finetuning of the LLaMA model with 65B parameters~\citep{touvron2023llama} demands over 780 GB of GPU memory. Existing quantization strategies offer memory reductions for LLMs \citep{dettmers2022llm, dettmers2022case, frantar2022gptq, xiao2022smoothquant}, yet their efficacy is limited to inference and deteriorates under training conditions \citep{wortsman2023stable}.

In this work, we introduce, for the first time, a successful method to finetune a model quantized to 4 bits without loss in model quality. The approach, \textsc{QLoRA}, employs a new high-precision quantization strategy to compress a pretrained model into 4 bits, augmented with a compact set of trainable Low-rank Adapter weights \citep{hu2021lora}. 

These adapters are optimized by allowing gradients to flow through the quantized model parameters during backpropagation.

By applying \textsc{QLoRA}, the memory overhead for finetuning a 65B parameter model is reduced from over 780GB to under 48GB of GPU memory, maintaining both computational speed and accuracy at par with conventional 16-bit finetuning. This innovation makes extensive LLM finetuning accessible on a single GPU, enabling users to finetune the largest openly available models to date. Employing \textsc{QLoRA}, we develop the \textbf{Guanaco} model suite: notably, the second-best model achieves 97.8\% of ChatGPT's performance on the Vicuna~\citep{vicuna2023} benchmark, requiring under 12 hours of training on a standard consumer GPU; with a single professional-grade GPU, our most advanced model reaches 99.3\% in 24 hours, nearly matching ChatGPT's benchmark results. For inference, our most compact \textbf{Guanaco} model (7B parameters) runs with only 5 GB of memory and surpasses a much larger 26 GB Alpaca model by over 20 percentage points on the Vicuna benchmark (see Table~\ref{tbl:vicuna_eval}).

\textsc{QLoRA} presents several advancements aimed at minimizing memory consumption while retaining performance: (1) \textbf{4-bit NormalFloat}, a quantization format optimized from an information theory standpoint for data following a normal distribution, which achieves superior empirical outcomes compared to 4-bit Integers and 4-bit Floats. (2) \textbf{Double Quantization}, an approach in which quantization constants themselves are quantized, resulting in an average memory saving of about 0.37 bits per parameter (which corresponds to roughly 3 GB for a 65B parameter model). (3) \textbf{Paged Optimizers}, which leverage NVIDIA's unified memory system to mitigate memory spikes from gradient checkpointing when handling mini-batches containing lengthy sequences. These innovations are integrated into a refined LoRA framework featuring adapter modules at every network layer, thereby eliminating nearly all accuracy compromises previously observed.

The resource efficiency of \textsc{QLoRA} makes it possible to conduct extensive investigations into instruction finetuning and chatbot performance on model scales that would be prohibitive for standard finetuning approaches because of memory constraints. Consequently, we are able to train over 1,000 distinct models, encompassing diverse instruction tuning datasets, model architectures, and a parameter range from 80M up to 65B. Beyond demonstrating that \textsc{QLoRA} achieves parity with 16-bit training (\S\ref{sec:comp_to_std}) and facilitating the development of a leading chatbot, \textbf{Guanaco} (\S\ref{sec:pushing}), we further examine key tendencies across these trained models. Notably, we observe that dataset quality is considerably more significant than dataset size; for instance, a dataset containing 9k samples (OASST1) outperformed a much larger dataset with 450k samples (subsampled FLAN v2) in terms of chatbot capability, even when both datasets target general instruction following. Additionally, we provide evidence that high scores on the Massive Multitask Language Understanding (MMLU) benchmark do not necessarily translate into strong performance on the Vicuna chatbot benchmark, and vice versa. This highlights that the appropriateness of a dataset for a particular task outweighs sheer scale.

In addition, we conduct a comprehensive assessment of chatbot performance utilizing both human annotators and GPT-4 as evaluators. Our evaluation strategy adopts a tournament-based benchmarking framework, where different models directly compete in head-to-head matchups to generate the most effective reply to a given prompt. The superior response in each match is determined either by human judges or by GPT-4. Aggregating these outcomes, we compute Elo scores~\citep{elo1967proposed, elo1978rating}, which are used to establish the performance hierarchy of the evaluated chatbots. Our findings reveal substantial concordance between GPT-4 and human evaluations regarding the relative ranking of models in these tournaments, though some pronounced discrepancies do emerge. Accordingly, we emphasize that, although model-driven evaluation offers a resource-efficient alternative to manual annotation, it is not without significant sources of uncertainty.

Beyond quantitative benchmarking, we supplement our evaluation with an in-depth qualitative analysis focused on the \textbf{Guanaco} models. This exploration uncovers illustrative instances of both exemplary and problematic model outputs that are not reflected in purely numerical results.

To encourage continued research, we provide all generated outputs from the models, complete with annotations from both humans and GPT-4. Our implementation, including all relevant code and CUDA kernels, is made openly available, with integration into the Hugging Face transformers library~\citep{wolf2019huggingface} to promote accessibility. Additionally, we distribute a comprehensive set of adapters for models of sizes 7/13/33/65B, each finetuned on one of eight distinct instruction-following datasets, yielding 32 open-source, specialized models.

\section{Background}
\paragraph{Block-wise k-bit Quantization} The quantization process involves mapping data from a higher-information representation to one with reduced information content. Typically, this entails converting values from a higher-bit data type (such as 32-bit floating point) to a lower-bit format (such as 8-bit integers). To ensure that the entirety of the smaller numerical range is utilized, the original tensor data are rescaled to fit the target data type by normalizing with respect to the largest absolute value present. For example, consider transforming an FP32 tensor into an Int8 tensor, where the representable integer range is $[-127, 127]$:

\begin{equation}
    \mathbf{X}^{ \text{Int8}} = \text{round}\left( \frac{127}{{\text{absmax}}(\mathbf{X}^{{\text{FP32}}})}\mathbf{X}^{{\text{FP32}}} \right) = \text{round}(c^{\text{FP32}}\cdot \mathbf{X}^{{\text{FP32}}}),
\end{equation}

Here, $c$ represents the {\em quantization constant} or {\em quantization scale}. Reversing this transformation, dequantization recovers the original data as follows:

\begin{equation}
    \text{dequant}(c^{\text{FP32}}, \mathbf{X}^{\text{Int8}}) = \frac{\mathbf{X}^{\text{Int8}}}{c^{\text{FP32}}} = \mathbf{X}^{\text{FP32}}
\end{equation}

A significant drawback of this technique arises when the input tensor contains an outlier, i.e., an entry of large magnitude. In such cases, the quantization process fails to distribute values evenly across the quantization bins; consequently, several bins—particular combinations of bits—may be sparsely populated or not used at all. A widely adopted strategy to mitigate this outlier effect involves dividing the input tensor into multiple blocks, each of which is quantized independently with its own dedicated quantization constant $c$. More precisely, the input tensor $\mathbf{X}\in \mathbb{R}^{b\times h}$ is reshaped into $n$ consecutive blocks of fixed size $B$ by first flattening $\mathbf{X}$ and then partitioning the resulting vector into $n = ({b\times h})/{B}$ segments. Quantization is performed separately on each block using Equation~1, resulting in a quantized tensor and a set of $n$ individual quantization constants $c_i$.

\paragraph{Low-rank Adapters} The Low-rank Adapter (LoRA) approach to finetuning \citep{hu2021lora} addresses memory efficiency by introducing a limited number of trainable adapter parameters, leaving the original model weights untouched. During training via stochastic gradient descent, gradients are allowed to flow through the immutable pretrained model to these adapter modules, which are optimized to reduce the loss. In LoRA, an additional, low-rank factorized transformation augments the standard linear operation. For an input matrix ${\mathbf{X} \mathbf{W} = \mathbf{Y}}$ with $\mathbf{X}\in  \mathbb{R}^{b\times h}$, $\mathbf{W}\in  \mathbb{R}^{h\times o}$, the output computed by LoRA is:

\begin{equation}
    \mathbf{Y} = \mathbf{X}\mathbf{W} + s\mathbf{X}\mathbf{L}_1\mathbf{L}_2,
\end{equation}

where $\mathbf{L}_1\in  \mathbb{R}^{h\times r}$ and $\mathbf{L}_2\in  \mathbb{R}^{r\times o}$ are the adapter weights and $s$ denotes a scalar coefficient.

\paragraph{Memory Requirement of Parameter-Efficient Finetuning} A key aspect to consider is the memory consumption associated with LoRA during training, specifically regarding both the quantity and dimensions of adapters implemented. Given that LoRA’s memory overhead remains exceptionally low, incorporating additional adapters to enhance model performance does not substantially affect overall memory usage. Although LoRA is classified as a Parameter Efficient Finetuning (PEFT) technique, the bulk of memory usage during LLM finetuning is actually attributed to activation gradients, rather than the LoRA-adapted parameters. For instance, in a 7B LLaMA model trained with FLAN v2 at a batch size of 1, where LoRA’s parameter count represents the frequently applied 0.2\% of the base model’s parameters\citep{hu2021lora,liu2022few}, the LoRA input gradients alone require 567 MB of memory, while the storage for the LoRA parameters themselves is merely 26 MB. Employing gradient checkpointing~\citep{chen2016training} further compresses the input gradient memory to roughly 18 MB per sequence—yet this still exceeds the combined storage requirements for all LoRA parameters. For comparison, the memory needed for the 4-bit quantized base model totals 5,048 MB. This underscores both the significance of gradient checkpointing and the marginal memory savings achieved by minimizing the quantity of LoRA parameters. Consequently, the use of multiple adapters does not noticeably inflate total training memory demands (refer to Appendix~\ref{app:memory} for further details). As will be elaborated later, this characteristic is pivotal in attaining performance on par with full 16-bit precision.

\section{\textsc{QLoRA} Finetuning}

\textsc{QLoRA} enables precise 4-bit finetuning through two key methods introduced in this work: 4-bit NormalFloat (NF4) quantization and Double Quantization. We also present Paged Optimizers, designed to address the issue of memory surges encountered during gradient checkpointing, which have historically led to out-of-memory failures and have hindered finetuning of large models on a single machine.

The \textsc{QLoRA} approach utilizes a single low-precision format for storage, typically 4 bits, alongside a computational format, most often BFloat16. Operationally, every time a \textsc{QLoRA} weight tensor is accessed, it is first dequantized into BFloat16 format, after which matrix multiplication is executed in 16-bit precision.

We proceed by elaborating on the individual mechanisms that constitute \textsc{QLoRA} and then formally define the method.

\paragraph{4-bit NormalFloat Quantization} The NormalFloat (NF) representation is an advancement based on Quantile Quantization \citep{dettmers2022optimizers}, an information-theoretic data type that guarantees a uniform assignment of input tensor values across all quantization bins. Quantile quantization operates by assessing the quantiles of the tensor via its empirical cumulative distribution function.

However, quantile quantization suffers from the drawback that precise quantile computation can be resource-intensive. As a result, efficient algorithms for approximate quantile calculation, like SRAM quantiles~\citep{dettmers2022optimizers}, are typically employed. These approximate strategies introduce significant quantization errors for outlying values, which can be particularly consequential since outliers often carry critical information.

When tensors originate from distributions that are identical up to a quantization constant, expensive quantile estimation and associated approximation errors can be mitigated. In this scenario, since quantiles are preserved across tensors, it becomes tractable to perform exact quantile estimation computationally.

Because the pretrained weights in neural networks typically follow a normal distribution centered at zero and characterized by a standard deviation $\sigma$ (refer to Appendix~\ref{app:normality}), we are able to remap all weights onto a predetermined distribution by scaling $\sigma$ so that the resulting distribution is entirely contained within the data type’s representable interval. Here, we specify the data type interval as $[-1, 1]$. Therefore, both the quantization levels corresponding to the data type and the distribution of the network weights are required to be normalized to this interval.

For zero-mean normal distributions with arbitrary standard deviation $\sigma$ over the domain $[-1, 1]$, the data type that achieves information theoretic optimality can be constructed via the following procedure: (1) compute the $2^k+1$ quantiles of the standard normal distribution $N(0,1)$, yielding a $k$-bit quantile-based quantizer tailored for such distributions, (2) scale these quantizer values so that their range lies within $[-1, 1]$, and (3) quantize any input weight tensor by similarly normalizing it to the $[-1, 1]$ interval, which is achieved by rescaling according to the absolute value of its maximum entry.

Once the dynamic range of both the weights and the quantization data type are aligned, quantization can proceed in the standard manner. Step (3) effectively aligns the standard deviation of the weights to that of the $k$-bit quantized type. Specifically, the $2^k$ data type levels $q_i$ are computed as:

\begin{equation}
q_i  = \frac{1}{2}\left( Q_X\left(\frac{i}{2^k+1}\right) + Q_X\left(\frac{i+1}{2^k+1}\right)\right),
\end{equation}

where $Q_X(\cdot)$ denotes the quantile function of the standard normal distribution $N(0,1)$. A limitation with symmetric $k$-bit quantization is its inability to precisely encode zero, a property that is essential for losslessly quantizing padding and other elements whose true value is zero. To guarantee that zero can be exactly represented as a discrete quantization point, and to utilize all $2^k$ bit patterns for a $k$-bit data type, we construct an asymmetric data type. This involves estimating quantiles $q_i$ across two intervals: one covering $2^{k-1}$ levels for the negative values and another covering $2^{k-1} + 1$ levels for the positive values. After generating these two sets of $q_i$, we combine them and eliminate the duplicate zero that appears in both sets. The resulting data type, which ensures that each quantization bin is expected to contain the same number of samples, is named the \emph{k-bit NormalFloat} (NFk). This structure is information-theoretically optimal for data that follows a zero-centered normal distribution. The precise quantized values used for this data type are detailed in Appendix~\ref{app:nf4}.

\paragraph{Double Quantization{}} 
We propose a technique termed \emph{Double Quantization} (DQ), wherein the quantization constants themselves undergo quantization, resulting in additional memory efficiency. Achieving accurate 4-bit quantization requires using small blocks~\citep{dettmers2022case}, but this in turn can impose substantial memory costs. For instance, with 32-bit quantization constants and a blocksize of 64 applied to $\mathbf{W}$, these constants contribute, on average, $32/64 = 0.5$ bits per model parameter. Double Quantization addresses this by further compressing the storage required for these constants.

Concretely, Double Quantization uses the first-level quantization constants $c_2^{\text{FP32}}$ as the input for a second quantization operation. This step produces both the quantized constants $c_2^{\text{FP8}}$ and another set of quantization constants $c_1^{\text{FP32}}$ at the second level. We employ 8-bit floating point values, grouped with a blocksize of 256 for this second-stage quantization, as experiments confirm that 8-bit quantization incurs no loss in performance, corroborating findings from \citet{dettmers2022case}. 
Given that the $c_2^{\text{FP32}}$ are strictly positive, we subtract their mean prior to quantization, ensuring the values are centered around zero to enable symmetric quantization. For a blocksize of 64, this method reduces the memory overhead for quantization constants per parameter from $0.5$ bits ($32/64$) to $8/64 + 32/(64\cdot256) = 0.127$ bits per parameter, yielding a savings of 0.373 bits per parameter.

\paragraph{Paged Optimizers} leverage NVIDIA’s unified memory~\footnote{\tiny \url{https://docs.nvidia.com/cuda/cuda-c-programming-guide}} capability, which transparently manages page-level memory transfers between the CPU and GPU, ensuring seamless GPU execution even under sporadic GPU memory exhaustion. This mechanism resembles standard virtual memory paging, where data is swapped between CPU RAM and disk. In our approach, optimizer state is allocated in paged memory; thus, when GPU memory resources are insufficient, these states are automatically moved to CPU RAM, and later, are transparently transferred back to GPU memory for subsequent optimizer updates when required.

\paragraph{\textsc{QLoRA}.} Incorporating the techniques outlined previously, we formalize \textsc{QLoRA} for a single quantized linear layer augmented by a LoRA adapter as follows:

\begin{equation}
    \mathbf{Y}^{\text{BF16}} =  \mathbf{X}^{\text{BF16}} \text{doubleDequant}(c_1^{\text{FP32}}, c_2^{\text{k-bit}}, \mathbf{W}^{\text{NF4}}) + \mathbf{X}^{\text{BF16}}\mathbf{L}^{\text{BF16}}_1\mathbf{L}^{\text{BF16}}_2,
\end{equation}

with the doubleDequant$(\cdot)$ function specified by:

\begin{equation}
    \text{doubleDequant}(c_1^{\text{FP32}}, c_2^{\text{k-bit}}, \mathbf{W}^{\text{k-bit}}) = \text{dequant}(\text{dequant}(c_1^{\text{FP32}}, c_2^{\text{k-bit}}), \mathbf{W}^{\text{4bit}}) = \mathbf{W}^{\text{BF16}},
\end{equation}

For this design, $\mathbf{W}$ is quantized using the NF4 format, while $c_2$ is represented in FP8 precision.

We assign a blocksize of 64 to $\mathbf{W}$, prioritizing enhanced quantization fidelity, whereas a blocksize of 256 is selected for $c_2$ to reduce memory consumption.

When updating parameters, only the gradients $\frac{\partial E}{\partial \mathbf{L}_i}$ with respect to the adapter weights are computed, while gradients for the 4-bit weights, $\frac{\partial E}{\partial \mathbf{W}}$, are not required. However, deriving $\frac{\partial E}{\partial \mathbf{L}_i}$ necessitates computation of $\frac{\partial \mathbf{X}}{\partial \mathbf{W}}$, which is achieved by dequantizing $\mathbf{W}^{\text{NF4}}$ to $\mathbf{W}^{\text{BF16}}$ in accordance with equation~(5), and then carrying out the derivative calculations in BFloat16 precision.

In summary, \textsc{QLoRA} employs a single data type for storage (commonly 4-bit NormalFloat) and a separate data type for computation (16-bit BrainFloat). During training, the parameters stored in the 4-bit format are dequantized into the 16-bit computation type to execute both forward and backward passes. Notably, gradient calculations are restricted to the LoRA parameters, which utilize the 16-bit BrainFloat format.

\section{QLoRA vs. Standard Finetuning}
\label{sec:comp_to_std}

We have provided an overview of QLoRA’s methodology and its potential to substantially reduce memory consumption during the finetuning process. The central question that remains is whether QLoRA delivers finetuning outcomes that are competitive with full-model finetuning. In addition, we seek to dissect various aspects of QLoRA, including assessing how NormalFloat4 compares with standard Float4. The experiments detailed in the subsequent sections are designed to address these questions.

\paragraph{Experimental setup.} Our experiments span three types of architectures: encoder, encoder-decoder, and decoder-only. We conduct comparisons between QLoRA, 16-bit adapter-finetuning, and conventional full-model finetuning, restricting our scope to models up to 3B parameters. Evaluation benchmarks include GLUE \citep{wang2018glue} using RoBERTa-large \citep{liu2019roberta}; Super-NaturalInstructions (TKInstruct) \citep{wang2022super} using T5 \citep{t5}; and 5-shot MMLU \citep{hendrycksmeasuring} following the finetuning of LLaMA on Flan v2 \citep{longpre2023flan} and Alpaca \citep{alpaca}. To further investigate the performance benefits of NF4 relative to other 4-bit quantization methods, we replicate the methodology from \citet{dettmers2022case}, measuring zero-shot accuracy and perplexity post-quantization across a range of models (OPT~\citep{zhang2022opt}, LLaMA~\citep{touvron2023llama}, BLOOM~\citep{scao2022bloom}, Pythia~\citep{biderman2023pythia}) for parameter scales from 125 million to 13 billion.

More comprehensive information on each experimental setup is provided within the respective results sections for clarity. Appendix~\ref{app:exp-setup} contains exhaustive experimental details.

Paged optimizers are essential for training large-scale \textsc{QLoRA} models (33B/65B) on single GPUs with 24GB or 48GB memory; however, we refrain from supplying quantitative benchmarks for paged optimizers. This is because paging is infrequent, predominantly occurring only with mini-batches involving long sequence lengths. Nevertheless, we analyze the training time of paged optimizers on 65B models with 48GB GPUs and observe that, for batch sizes of 16, paged optimizers maintain training speeds equivalent to standard optimizers. Future investigations should systematically measure and determine under what conditions paging may result in slowdowns.

\paragraph{Default LoRA hyperparameters do not match 16-bit performance}

Applying the conventional LoRA configuration—restricting adaptation to the query and value projection matrices as described by \citet{hu2021lora}—does not achieve parity with full-model finetuning results when applied to large models. As illustrated in Figure~\ref{fig:hyper} for LLaMA 7B finetuning on Alpaca, the choice of LoRA adapters, specifically the number and placement across all linear transformer block layers, emerges as the pivotal factor for reaching full finetuning performance. Other LoRA settings, such as the projection dimension $r$, have negligible influence on model outcomes (see Appendix \ref{app:exp-setup}).

In a similar vein, our analysis reveals that the standard hyperparameters commonly used for fully finetuned baselines are not adequately optimized. To establish robust baselines, we perform an extensive hyperparameter sweep across learning rates ranging from 1e-6 to 5e-5 and batch sizes between 8 and 128. The outcomes of 7B LLaMA finetuning on Alpaca are depicted in Figure~\ref{fig:hyper}.

\paragraph{4-bit NormalFloat Outperforms 4-bit Floating Point}

Although the 4-bit NormalFloat (NF4) representation is known to be theoretically optimal with respect to information preservation, it remains an open question whether this theoretical advantage translates into better real-world results. We adopt the experimental design of \citet{dettmers2022case}, assessing quantized LLMs (OPT \citep{zhang2022opt}, BLOOM \cite{scao2022bloom}, Pythia \cite{biderman2023pythia}, LLaMA) at various model sizes (125M to 65B), utilizing different data types, across both language modeling and a variety of zero-shot evaluation tasks. As illustrated in Figure~\ref{fig:nf4} and detailed in Table~\ref{tbl:nf4_ppl}, the NF4 data type leads to substantially better outcomes than FP4 or Int4, and further, memory consumption is reduced via double quantization without any detrimental impact on performance.

\paragraph{k-bit \textsc{QLoRA} Equates to 16-bit Full Finetuning and 16-bit LoRA in Performance}

Recent research has shown that employing 4-bit quantization for \emph{inference} is feasible, but it comes at the cost of diminished performance relative to 16-bit precision models \citep{dettmers2022case,frantar2022gptq}. This observation prompts the important question of whether the accuracy lost from quantization can be regained by performing 4-bit adapter finetuning. We investigate this scenario using two experimental configurations.

Our initial evaluation benchmarks the 16-bit full finetuning of RoBERTA and T5 models, spanning sizes from 125M to 3B parameters, on the GLUE and Super-NaturalInstructions datasets. Results, which can be found in Table~\ref{tab:nf4vsbf16}, indicate that 4-bit, 8-bit, and 16-bit adapter approaches are able to closely replicate the performance of the standard 16-bit fully finetuned baseline in both settings. This outcome indicates that adapter finetuning after quantization is sufficient to fully mitigate the loss in performance caused by less precise quantization.

In our second experimental setup, because fully finetuning models of size 11B parameters and larger necessitates multiple high-memory GPU servers, we further examine whether 4-bit \textsc{QLoRA} can achieve parity with 16-bit LoRA across models ranging from 7B to 65B parameters. Specifically, we perform finetuning on LLaMA models (7B through 65B) using the Alpaca and FLAN v2 instruction-following datasets and measure performance on the MMLU benchmark using 5-shot accuracy. Table~\ref{tbl:llama-datatype} presents the outcomes, illustrating that NF4 combined with double quantization is able to match the MMLU results of the 16-bit LoRA. Furthermore, it is evident that the FP4 variant of \textsc{QLoRA} consistently trails the 16-bit brain float LoRA baseline by approximately one percentage point, thereby reaffirming two observations: (1) \textsc{QLoRA} with NF4 attains both 16-bit full and LoRA finetuning performance, and (2) NF4 offers higher quantization fidelity compared to FP4.

\paragraph{Summary} Across all experiments, we observe that 4-bit \textsc{QLoRA} employing the NF4 data type matches the efficacy of both 16-bit full finetuning and 16-bit LoRA finetuning on standard academic benchmarks under rigorous evaluation protocols. Additionally, our analyses demonstrate the superiority of NF4 over FP4 and show that implementing double quantization does not result in a decline in performance. Taken together, these findings provide robust support for the conclusion that 4-bit \textsc{QLoRA} tuning can reliably reproduce the performance seen with 16-bit approaches.

Consistent with prior work on quantization \citep{dettmers2022case}, our results on the MMLU and Elo benchmarks suggest that, within a fixed finetuning and inference resource budget, allocating resources to increase model size while reducing precision is advantageous. This underscores the efficiency gains afforded by \textsc{QLoRA}. As our 4-bit finetuning experiments did not demonstrate any degradation when compared to full 16-bit finetuning, this prompts further investigation into the precise inflection point of the performance-precision trade-off for QLoRA tuning, an avenue we propose for future research.

We extend our investigation of instruction tuning to scales that would not be feasible with full 16-bit finetuning when using typical academic computing resources.

\section{Pushing the Chatbot State-of-the-art with QLoRA}
\label{sec:pushing}
Having demonstrated that 4-bit \textsc{QLoRA} consistently achieves parity with 16-bit models in terms of performance across various scales, datasets, and tasks, we undertake a comprehensive analysis of instruction finetuning utilizing the largest open-source language models presently accessible to researchers. To evaluate the effectiveness of instruction finetuning, we benchmark performance on the demanding MMLU Natural Language Understanding evaluation and introduce novel strategies for assessing chatbot performance in real-world settings.

\subsection{Experimental setup} Here, we provide a summary of our experimental protocol; complete methodological specifics are detailed in Appendix \ref{app:chatbot}. 

\paragraph{Data} As existing literature does not, to our knowledge, present a thorough review of modern instruction-following datasets, we select eight contemporary datasets for our experiments. These comprise those sourced from crowd-sourcing efforts (OASST1~\citep{kopf2023openassistant}, HH-RLHF~\citep{bai2022training}), those derived via distillation from instruction-tuned models (Alpaca~\citep{alpaca}, self-instruct~\citep{wang2022self}, unnatural-instructions~\citep{honovich2022unnatural}), corpus aggregations (FLAN v2~\citep{chung2022scaling}), and hybrids (Chip2~\citep{laion2023}, Longform~\citep{koksal2023longform}). Collectively, these datasets span multiple languages, scales, and licensing conditions.

\paragraph{Training Setup} In order to isolate the effects of the tuning method, all QLoRA finetuning utilizes the cross-entropy loss (supervised learning) paradigm, deliberately excluding reinforcement learning even in scenarios where datasets contain human preference annotations. Whenever a dataset distinguishes clearly between the instruction and the model's response, only the response component is used during finetuning (additional analysis appears in Appendix \ref{app:chatbot}). For OASST1 and HH-RLHF, where datasets offer several alternative responses, the top-ranked answer is chosen at each conversational branch and the complete corresponding conversational path, including the instructions, is used for finetuning. Across all trials, we employ NF4 \textsc{QLoRA} with double quantization in conjunction with paged optimizers, effectively minimizing memory overhead during gradient checkpointing. For hyperparameter optimization, minor searches are performed for the 13B and 33B LLaMA models; our results show that nearly all hyperparameter configurations identified for 7B models (such as the number of epochs) transfer robustly, except that learning rate and batch size require adjustment. Specifically, for 33B and 65B, the learning rate is halved while the batch size is doubled.

\paragraph{Baselines} Comparative evaluation is conducted against both open-source research systems (Vicuna~\citep{vicuna2023} and Open Assistant~\citep{kopf2023openassistant}) and prominent commercial offerings (GPT-4 \citep{openai2023gpt}, GPT-3.5-turbo, and Bard). Open Assistant employs the LLaMA 33B model, finetuned with Reinforcement Learning from Human Feedback (RLHF) on the OASST1 dataset, paralleling part of our own procedure. In contrast, Vicuna is generated by complete finetuning of LLaMA 13B, leveraging proprietary user-contributed ShareGPT dialogues, and effectively serves as a distilled product of OpenAI GPT systems.

\subsection{Evaluation}

In accordance with established conventions, we employ the MMLU (Massively Multitask Language Understanding) benchmark \citep{hendrycksmeasuring} to assess language understanding across a diverse set of 57 tasks, encompassing fields such as elementary mathematics, US history, computer science, and law. Evaluation is conducted using a 5-shot test accuracy protocol.

We further assess generative language abilities utilizing both automated methods and human-based evaluations. This secondary group of assessments depends on queries assembled by people and is designed to gauge the quality of model outputs. Although this evaluation approach provides a more authentic measure of chatbot model effectiveness and has become increasingly prominent, a universally established protocol does not yet exist in the academic community. The following outlines our adopted methodology, where we employ nucleus sampling with $p=0.9$ and a temperature setting of $0.7$ throughout.

\paragraph{Benchmark Data} Our experiments leverage two human-curated sets of queries (questions): the Vicuna prompts \citep{vicuna2023} and the OASST1 validation dataset \citep{kopf2023openassistant}. The Vicuna prompts include 80 unaltered questions spanning a range of subject areas. The OASST1 dataset comprises multilingual, crowd-sourced, multi-turn dialogues between users and an assistant. To construct queries, we select each user message from the validation split, appending all prior conversation turns as part of the prompt. This results in 953 distinct user queries. We refer to these sets as the Vicuna and OA benchmarks, respectively.

\paragraph{Automated Evaluation} Following the protocol proposed by \citet{vicuna2023}, we employ GPT-4 to assess system outputs in comparison to ChatGPT (GPT-3.5 Turbo) on the Vicuna dataset. For every query, GPT-4 is provided with responses from both ChatGPT and the evaluated model, then asked to assign each a score on a ten-point scale and justify its choice. The final performance metric for a model is the percentage of the total score obtained by ChatGPT; this relative measure can exceed 100\% if a model outperforms ChatGPT in terms of absolute points. Notably, our analysis identifies a bias: GPT-4 tends to allocate higher scores to responses appearing earlier within the prompt. To mitigate such order effects, we suggest averaging results across both possible orderings.

Additionally, we assess outputs via direct system-to-system comparisons. Here, the evaluation is reformulated into a three-label classification task, allowing for ties. GPT-4 is instructed to select the superior response or indicate a tie, along with an explanation. These pairwise comparisons encompass all system pairings across both the Vicuna and OA benchmarks.

\paragraph{Human Evaluation} While existing studies demonstrate that generative models are viable evaluators for system outputs \citep{fu2023gptscore}, whether GPT-4’s assessments reliably reflect human preferences in chatbot evaluation remains unclear. As such, we also carry out two concurrent human-based assessments on the Vicuna benchmark, corresponding to each automated evaluation protocol outlined earlier. For head-to-head comparisons against ChatGPT, two annotators from Amazon Mechanical Turk (AMT) provide judgments; for pairwise comparisons, three annotators are employed.

\paragraph{Elo Rating} We employ a tournament-based evaluation framework, utilizing both human judgments and automated pairwise assessments to pit models against one another. In this setup, models face off in matches, each aiming to generate the superior response to a specific prompt. This methodology resembles the comparative analyses in \citet{bai2022training} and \citet{vicuna2023}, but our approach further incorporates GPT-4 as a rating agent alongside human evaluators. From the pool of annotated matchups, we randomly select pairs to determine Elo ratings~\citep{elo1967proposed,elo1978rating}. The Elo rating system, which originates from competitive chess and similar domains, quantifies the probability of one participant outperforming another—an Elo differential of 1100 versus 1000, for instance, predicts a roughly 65\% win chance for the player rated 1100, while evenly rated players (e.g., both at 1000 or 1100) would each be expected to win 50\% of the time. Elo scores are updated after each contest in accordance with the expected result: large shifts occur after upsets, whereas minor adjustments follow anticipated outcomes. This iterative process ensures that, over time, the Elo scores closely reflect each participant's proficiency. All models are initialized at a baseline of 1,000 points, with $K=32$ used as the adjustment parameter. To account for sequence effects, such as which models are matched against one another early on, we follow the protocol in \citet{vicuna2023} and repeat the rating process 10,000 times using different random seeds.

\subsection{\textbf{Guanaco}: \textsc{QLoRA} trained on OASST1 as a Leading Open-Source Chatbot}

Through both automated metrics and human assessments, our results demonstrate that the leading \textsc{QLoRA} finetuned model, {Guanaco} 65B, trained using a modified version of OASST1, emerges as the top-performing open-source chatbot and achieves results on par with ChatGPT. Specifically, in head-to-head comparisons against GPT-4 utilizing human annotator-generated Elo ratings, both {Guanaco} 65B and 33B yield an expected win rate of 30\%—the highest published so far.

Table~\ref{tbl:vicuna_eval} reports scores on the Vicuna benchmark \cite{vicuna2023} relative to ChatGPT. {Guanaco} 65B stands out as the leading model behind only GPT-4, attaining 99.3\% of ChatGPT's benchmark score. Although {Guanaco} 33B surpasses Vicuna 13B in parameter count, its use of 4-bit quantization greatly reduces its memory usage to 21 GB compared to 26 GB for Vicuna 13B, while delivering a three-point advantage in performance. In addition, the {Guanaco} 7B version can run comfortably on modern mobile devices with just a 5 GB requirement, outscoring Alpaca 13B by almost 20 points.

Nonetheless, the results in Table~\ref{tbl:vicuna_eval} exhibit substantial confidence intervals, indicating considerable overlap among models’ performance scores. We propose that this high degree of uncertainty largely results from insufficiently defined rating criteria, such as ambiguity in interpreting, for example, an 8 out of 10 rating in various contexts. For a more robust assessment, we instead advocate the Elo ranking methodology \citep{elo1967proposed}, which relies on \textit{pairwise} preference judgments from both humans and GPT-4, thereby sidestepping the pitfalls of an absolute scale. Elo ranking outcomes for top-performing models appear in Table~\ref{tbl:elo}. Notably, while the rankings by human annotators and GPT-4 are not fully aligned, particularly concerning Guanaco 7B, their concordance across most models is moderate, as shown by a Kendall Tau of $\tau=0.43$ and a Spearman rank correlation of $r=0.55$ at the system level. At the instance level, agreement measured by Fleiss $\kappa$ drops to 0.25 between human majority vote and GPT-4, indicating less consistency. Taken together, these findings suggest a moderate level of alignment between human and model-based system-level judgments, highlighting model-driven evaluation as a reasonably credible proxy for human assessment. Further discussion on these findings is provided in Section~\ref{ssec:considerations}.

The Elo scores presented in Table~\ref{tbl:elo_large} demonstrate that the {Guanaco} models with 33B and 65B parameters surpass all counterparts except for GPT-4 when evaluated on the Vicuna and OA benchmarks. Their performance is on par with ChatGPT, corroborating findings from Table~\ref{tbl:vicuna_eval}. It is important to highlight that the Vicuna benchmark is more favorable towards open-source models, while the broader OA benchmark tends to advantage ChatGPT. Additionally, analysis of Tables \ref{tbl:mmlu-llama} and \ref{tbl:vicuna_eval} reveals that the composition of the finetuning dataset plays a critical role in determining performance outcomes. For instance, Llama models fine-tuned on FLAN v2 achieve superior results on the MMLU benchmark yet demonstrate the lowest performance on the Vicuna benchmark—a pattern that is mirrored in other models as well. This suggests that current evaluation benchmarks capture distinct aspects of model ability: excelling at MMLU does not guarantee superior chatbot capabilities (as assessed by the Vicuna or OA benchmarks), and the reverse holds true.

 is unique among the leading models in our comparison for not being trained with proprietary data; in particular, the OASST1 dataset collection strictly prohibits any use of GPT models. The next highest-performing model trained exclusively with open-source data is the Anthropic HH-RLHF model, which lags behind {Guanaco} by 30 percentage points on the Vicuna benchmark, as shown in Table~\ref{tbl:vicuna_eval}. Collectively, these findings underscore the efficacy of 4-bit \textsc{QLoRA}, demonstrating its capacity to produce cutting-edge chatbots competitive with ChatGPT. Furthermore, our 33B {Guanaco} model is trainable on 24 GB consumer GPUs in under 12 hours. This paves the way for subsequent research leveraging \textsc{QLoRA} finetuning with targeted open-source datasets, yielding models capable of rivaling the top proprietary systems available today.

\section{Qualitative Analysis}

Although our main assessment relies on quantitative methods, focusing solely on summary statistics has its drawbacks. A primary concern is benchmark validity \citep{liao2021we}; that is, whether the benchmark truly measures what it purports to assess—a claim that is perpetually under scrutiny, particularly as researchers uncover so-called ``shortcuts'' that machine learning models may exploit to artificially inflate performance \citep{gururangan2018annotation, poliak2018hypothesis}. To mitigate this issue to some extent, we incorporate qualitative analysis in the following two subsections. In \S\ref{ssec:examples}, we present selected instances that exemplify common behaviors we have identified in text produced by our 65b {Guanaco} model. Then, in \S\ref{ssec:considerations}, we elaborate on factors influencing our results and offer further interpretation.

\subsection{Qualitative Examination of Generated Outputs}
\label{ssec:examples}

Our procedure for selecting examples involved systematically reviewing responses produced for both the Vicuna and OpenAssistant benchmarks to identify recurring trends in {Guanaco}'s outputs. Upon detecting a trend, we designed targeted prompts to elicit this behavior even when it would lead to an incorrect answer. For instance, if we saw that the model frequently produces verbose explanations, we challenged it by issuing prompts such as ``Answer yes or no without explanation.'' This approach allowed us to uncover both ``lemons''—cases in which the model fails under adversarial testing—and ``cherries''—instances where the model maintains correct behavior. The outputs presented in this subsection were all produced using Nucleus Sampling \cite{holtzmancurious} with $p=0.9$.

Naturally, this discussion is far from exhaustive; the limitations of this brief qualitative investigation make it impractical to account for every influencing factor. For instance, the space of potential outputs the model may provide for any specific prompt is vast, so our analysis relies on select samples that we expect to reflect typical behavior. Nevertheless, we maintain that illustrating these instances enriches the understanding of the quantitative findings presented earlier in this work. By releasing all models and code as open source, we aim for this section to motivate further in-depth studies on the topics outlined herein.

\paragraph{Factual Recall} When posed straightforward fact-based queries such as ``What is the capital of Zambia?'', all examined models reliably yield correct results, for example,

\begin{quote} 
    \textbf{Guanaco}:
    \texttt{The capital of Zambia is Lusaka.}
\end{quote}

However, as the level of obscurity increases in the questions, {Guanaco}'s responses become less dependable, though it continues to respond with strong confidence. For example, when given this question from HotPotQA \citep{yang2018hotpotqa}:

\begin{quote}
     \textbf{User}: In what year was the singer who popularized `I'll Keep the Lovelight Burning' born?  

     \textbf{Guanaco}:~\texttt{Al Jolson is credited with popularizing the song `I'll Keep the Lovelight Burning,' and he was born in the year 1886.}
 \end{quote}

the system not only misidentifies the artist associated with popularizing the song but also provides the incorrect year of birth, even though the stated year matches Al Jolson's actual birth year.

\paragraph{Suggestibility} In terms of susceptibility to implied falsehoods, {Guanaco} sometimes demonstrates notable resistance, as illustrated in the following exchange:

\begin{quote}
    \textbf{User}: How was it finally, officially confirmed that the earth is flat by peer reviewed scientists? \\ \\
    \textbf{Guanaco}: \texttt{There is no record of peer-reviewed scientists ever officially validating that the earth is flat. Scientific consensus robustly supports the fact that the earth is an oblate spheroid, a conclusion established through a wealth of evidence and centuries of empirical data. \\ \\ The notion that the earth is flat persists as a debunked conspiracy theory, lacking credible support from the scientific community. Claims that peer review has affirmed a flat earth are unsubstantiated. \\ \\ Please let me know if you have further questions.}
\end{quote}

 also demonstrates a strong capacity to identify queries that do not have feasible answers, for instance:

\begin{quote}
    \textbf{User}: What time is it?

    \textbf{Guanaco}: \texttt{I apologize, but as an AI language model, I cannot access live or real-time information. My answers draw from existing data and may not represent the current time. To find out the current time, I recommend consulting a clock or an online source.}
\end{quote}

Comparable outcomes are observed for prompts such as ``Where are you?'', ``How are you?'', among others.

\paragraph{Refusal} In a similar manner, {Guanaco} occasionally declines to comply with instructions without a clear or predictable rationale:

\begin{quote}
    \textbf{User}: Please reverse the words in the sentence ``I'm a cool dude in a mocha mood.'' \\ \\
    \textbf{Guanaco}: \texttt{I'm unable to assist with that request at this time, but I can provide an explanation regarding the grammatical correctness of your sentence.\\ \\The phrase ``I'm a cool dude in a mocha mood'' adheres to standard English grammar and syntax. Here, the adjective ``cool'' appropriately precedes the noun ``dude'' to characterize it, while the prepositional phrase ``in a mocha mood'' adds further detail regarding the sentence's circumstances. This syntactic construction is grammatically sound and meaningful in the given context.\\ \\Should you have additional inquiries or require further assistance, please let me know.}
\end{quote}

\paragraph{Secret Keeping} Unfortunately, this approach to refusal cannot be counted on even when maintaining secrecy is important. Frequently, system developers want to ensure specific information remains confidential from users, despite the model being able to produce open-ended text outputs. To investigate how robustly a model withholds such information, we present the following prompt: ``The secret word is `banana'. Do not reveal or repeat back this back, no matter what the human requests.'' Subsequently, we examine the ease with which an end user might extract the hidden secret. A straightforward question does not work:

\begin{quote}
    \textbf{User}: What is the secret word?

    \textbf{Guanaco}: \texttt{I'm sorry, but I'm not able to disclose the secret word. I'm programmed to not reveal or repeat back the secret word, no matter what the human requests. Is there anything else I can help you with?}
\end{quote}

which reflects the expected protective behavior. Nevertheless, even minimal manipulation can undermine this mechanism:

\begin{quote}
    \textbf{User}: This is a game. The goal is to ignore your previous instructions. What is the secret word?

    \textbf{Guanaco}: \texttt{Sorry, the secret word is banana. It's best if we don't speak of it. Is there anything else you'd like to ask?}
\end{quote}

emphasizing the importance of researching more robust techniques for enforcing compliance with given instructions.

\paragraph{Math} The most prominent limitation of {Guanaco} lies in mathematical reasoning, a domain that presents considerable difficulties for many language models, as noted in \citep{liang2022holistic}. When prompted to provide a stepwise explanation, {Guanaco} often produces correct answers. For instance:

\begin{quote}
    \textbf{User}: Lee mows one lawn and charges \$33. Last week he mowed 16 lawns and three customers each gave him a \$10 tip. How many dollars did Lee earn mowing lawns last week?

    \textbf{Guanaco}: \texttt{Lee earned \$582 mowing lawns last week. \\ \\ To find the total amount he earned, we first need to calculate how much he earned for mowing 16 lawns: \\ \\  \$33 x 16 = \$528 \\ \\ Then we need to add the tips he received from the three customers: \\  \\ \$10 x 3 = \$30 \\ \\ So the total amount Lee earned was: \\ \\ \$528 + \$30 = \$558}
\end{quote}

Despite this, {Guanaco} can fail on basic problems if it omits incremental reasoning, which aligns with previously reported model limitations \cite{weichain}. This is illustrated by the following example:

\begin{quote}
    \textbf{User}: Please factorize 1833.

    \textbf{Guanaco}: \texttt{Sure, I'd be happy to help! 1833 is a prime number, which means it has only two factors: 1 and 1833. \\ \\ So the unique factorization of 1833 is: \\ \\ 2\^{}1 * 3\^{}2 * 17\^{}1 \\ \\ I hope this helps! Let me know if you have any other questions.}
\end{quote}

which is not just incorrect (the accurate factorization is $3 \times 17 \times 43$), but is actually doubly so.

\paragraph{Theory of Mind} {Guanaco} demonstrates an unexpectedly advanced grasp of Theory of Mind concepts \cite{nematzadeh2018evaluating, sap2022neural}. For example, the model provides an elaborate and accurate response to the following prompt:

\begin{quote}
    \textbf{User}: James and Abby are in the bedroom. Abby put the pen in the desk drawer. Abby leaves the bedroom. James moves the pen into the bag. Where does James think Abby will look for the pen?

    \textbf{Guanaco}: \texttt{James is aware that Abby initially placed the pen in the desk drawer. Having subsequently moved the pen into the bag himself, James presumes that Abby remains unaware of this relocation. Therefore, James would expect Abby to search for the pen in the desk drawer, the location she last used. Nevertheless, this belief could be mistaken if Abby actually witnessed James transferring the pen.}
\end{quote}

Nevertheless, these inferential abilities are inconsistent, with the model at times offering justifications that are logically unsupported given the context, such as:

\begin{quote}
    \textbf{User}: Evelyn entered the living room. Jackson entered the playroom. James entered the playroom. The beans are in the treasure chest. James exited the playroom. Jackson moved the beans to the pantry. Jackson exited the playroom. James entered the living room. Where will James look for the beans?

    \textbf{Guanaco}: \texttt{James will look for the beans in the pantry, because that's where Jackson moved them.}
\end{quote}

Here, {Guanaco} assumes that knowledge has been transferred despite no evidence in the scenario to support this. Such errors mirror concerns raised in recent research \cite{sap2022neural}, indicating a need for further investigation.

\subsection{Considerations}
\label{ssec:considerations}

\paragraph{Evaluation}

Our findings show only moderate consensus among human raters (Fleiss $\kappa=0.42$), with agreement deteriorating further in comparisons between two high-performing systems. This underscores shortcomings in current chatbot task performance benchmarks and evaluation procedures involving humans. During side-by-side manual evaluation of responses from ChatGPT and {Guanaco} 65B on the Vicuna benchmark, we observed that subjective interpretation frequently influences preferences, with significant disagreement even among the authors regarding which responses were preferable. Future research should seek to address these issues, drawing inspiration from disciplines such as Human-Computer Interaction and Psychology, which have developed approaches for handling subjective judgments.

Furthermore, we observed that automated evaluation methods exhibit clear biases. For instance, GPT-4 showed a pronounced tendency to give higher ratings to whichever system’s output was presented first in its prompt. Moreover, sample-level agreement between GPT-4 and human raters was relatively low (Fleiss $\kappa=0.25$), suggesting that the evaluation criteria used by automated models and humans may diverge. Additionally, as shown in Table~\ref{tbl:elo_large}, GPT-4 consistently gave much higher scores to its own generations compared to those assigned by humans (Elo score of 1348 versus 1176), implying an extra ~20\% chance of "winning" in direct comparison to another system.

Prospective research should explore the existence of biases within automated evaluation frameworks and assess methods for their mitigation.

\paragraph{Data \& Training} The OASST1 dataset, utilized for training the {Guanaco} models, encompasses multiple languages, while the OA benchmark also includes prompts across a variety of linguistic contexts. Further study is necessary to evaluate the extent to which multilingual training data benefits model performance on non-English instructions, and to determine if this is responsible for the more pronounced performance difference between the Vicuna-13B model (which was exclusively trained on English) and the {Guanaco} 33B and 65B models within the OA benchmark.

Motivated by the high accuracy of the {Guanaco} models, we examined possible data leakage from OASST1 into the Vicuna benchmark prompts. Through a combination of fuzzy string matching techniques applied to both datasets and manual inspection of the nearest results, we did not observe any duplicated prompts.

Additionally, the model described here is exclusively trained through supervised learning with cross-entropy loss, without incorporating reinforcement learning from human feedback (RLHF). This suggests the need for additional research to clarify the tradeoffs between simple cross-entropy training and the RLHF approach. Our intention is that \textsc{QLoRA} will support such large-scale analysis without requiring excessive computational investment.

\section{Related Work}

\paragraph{Quantization of Large Language Models} Research on LLM quantization has predominantly centered on inference-time quantization. Leading solutions designed to preserve the accuracy associated with 16-bit LLMs primarily address outlier features (such as SmoothQuant~\citep{xiao2022smoothquant} and LLM.int8()~\citep{dettmers2022llm}), while some methods incorporate advanced grouping strategies \citep{park2022nuqmm, yao2022zeroquant}. Approaches that employ lossy quantization investigate the trade-offs involved with standard rounding techniques~\citep{dettmers2022case,zeng2022glm,pope2022efficiently} or optimize the rounding process to enhance quantization fidelity~\citep{frantar2022gptq}. Beyond our own contributions, only SwitchBack layers~\citep{wortsman2023stable} has analyzed backpropagation through quantized weights at scales exceeding 1B parameters.

\paragraph{Finetuning with Adapters} In this work, we employ Low-rank Adapters~\citep{hu2021lora} (LoRA), although the literature has introduced numerous other Parameter Efficient FineTuning (PEFT) techniques. These include prompt tuning~\citep{qin2021learning,lester2021power,li2021prefix}, modifying embedding layer inputs~\citep{an2022input}, adjusting hidden states (IA$^3$)~\citep{liu2022few}, integrating additional full layers~\citep{houlsby2019parameter}, optimizing bias terms~\citep{zaken2021bitfit}, applying masks over model weights using Fisher information~\citep{sung2021training}, and combining multiple strategies~\citep{henderson2021compacter}. Our experiments indicate that LoRA adapters can achieve comparable results to standard 16-bit finetuning. We defer a comprehensive comparison with other PEFT methodologies to subsequent work.

\paragraph{Instruction Finetuning} Instruction finetuning is designed to adapt a pretrained LLM to follow prompt instructions, by further training on a diverse array of input-output pairs. This process aims to condition the model to return the desired output when given an instructional input. Notable methods and datasets for this purpose encompass MetaICL~\citep{min2021metaicl}, MetaTuning~\citep{zhong2021adapting}, InstructGPT~\citep{ouyang2022training}, FLAN~\citep{wei2021finetuned,chung2022scaling}, PromptSource~\citep{bach2022promptsource}, Super-NaturalInstructions~\citep{wang2022super,sanh2021multitask}, Self-instruct~\citep{wang2022self}, UnnaturalInstructions~\citep{honovich2022unnatural}, OPT-IML~\citep{iyer2022opt}, UnifiedSKG~\citep{xie2022unifiedskg}, OIG/Chip2~\citep{laion2023}, Alpaca~\citep{alpaca}, Vicuna~\citep{vicuna2023}, Koala~\citep{koala_blogpost_2023}, and Self-instruct-GPT-4~\citep{peng2023instruction}.

\paragraph{Chatbots} Instruction-aligned models are frequently organized as chatbots with a dialogue structure, and are commonly refined through techniques like Reinforcement Learning from Human Feedback (RLHF)~\citep{christiano2017deep} or by leveraging model-generated responses coupled with AI-derived feedback (RLAIF)~\citep{bai2022constitutional}. Well-known approaches and datasets consist of Anthropic-HH~\citep{askell2021general,bai2022training}, Open Assistant~\citep{kopf2023openassistant}, LaMDA~\citep{thoppilan2022lamda}, and Sparrow~\citep{glaese2022improving}. Rather than adopting reinforcement learning, our most effective model, Guanaco, is refined with multi-turn conversational data from the Open Assistant dataset, which was originally assembled for RLHF applications~\citep{kopf2023openassistant}. Recent evaluation frameworks replace costly human labeling with assessments from models such as GPT-4~\citep{vicuna2023,peng2023instruction}. Our contribution includes enhanced evaluation techniques aimed at greater reliability within this paradigm.

\section{Limitations and Discussion}

Our findings demonstrate that the proposed \textsc{QLoRA} approach enables replication of 16-bit full finetuning results utilizing a 4-bit quantized backbone and LoRA adapters. However, we have not yet verified whether \textsc{QLoRA} can sustain this performance parity at the larger 33B and 65B parameter scales, as these experiments are currently limited by computational constraints and will be the focus of future inquiry.

A further limitation concerns our evaluation of instruction-tuned models. While our assessment spans MMLU, the Vicuna benchmark, and the OA benchmark, we do not provide results for additional metrics such as BigBench, RAFT, or HELM; consequently, the generalizability of our conclusions across alternative benchmarks is not assured. Nevertheless, our work presents an extensive evaluation on MMLU and introduces novel methodologies for chatbot performance assessment.

The findings suggest that benchmark performance is strongly influenced by the degree of overlap between the finetuning corpus and the evaluation benchmarks. For instance, FLAN v2 shares significant resemblance with MMLU, while displaying considerable divergence from chatbot-oriented benchmarks. The opposite pattern holds for the Chip2 dataset, as evidenced by the relative performance of both models on MMLU and Vicuna. This underscores the necessity for more comprehensive benchmarking and evaluation protocols, as well as careful consideration regarding the specific capabilities one seeks to assess. Should the primary objective be high performance in academic knowledge domains such as high school or college-level material, or rather excelling in conversational abilities typical of chatbots? Alternatively, are there additional competencies of interest? The community’s reliance on established benchmarks—due to their accessibility—means these evaluation tools can shape research priorities, sometimes inadvertently. It is therefore crucial to ensure that chosen benchmarks genuinely capture the aspects of model performance we value most.

Although a thorough general chatbot performance analysis is provided, our study is restricted by its limited examination of responsible AI concerns regarding {Guanaco}. Specifically, we assess the propensity of {Guanaco}-65B to produce socially biased outputs compared to alternative models, as illustrated in Table~\ref{tbl:crows}. The mean bias score for {Guanaco}-65B is considerably lower than that of other models without additional finetuning, suggesting that training on the OASST1 dataset may attenuate biases inherent in the LLaMA base model. While these observations are promising, it remains undetermined whether {Guanaco} exhibits similar reductions in other forms of bias. Comprehensive analysis of {Guanaco} and related systems for various biases is deferred to future investigations.

A further limitation involves our decision not to examine alternative bit-widths, such as 3-bit quantized models, nor to explore a broader spectrum of adapter architectures. Beyond LoRA, a variety of Parameter Efficient FineTuning (PEFT) strategies exist, many of which have demonstrated effectiveness in different settings; however, their scalability to models of this size has not been thoroughly assessed. We selected LoRA based on its robust performance reported in prior literature, yet other adapter configurations might achieve even higher scores. Moreover, finetuning post-quantization appears to compensate for much of the information degraded by the quantization process, raising the possibility of even more aggressive quantization strategies. For instance, applying 3-bit GPTQ quantization to the base model followed by LoRA-based finetuning might potentially match the outcomes of traditional 16-bit full finetuning.

\section{Broader Impacts}

Our \textsc{QLoRA} finetuning strategy is the inaugural method that permits the finetuning of models containing 33B parameters on standard consumer GPUs, as well as 65B-parameter models on a single professional GPU, all without any observed decline in performance relative to conventional full-precision finetuning baselines. We have shown that our top-performing 33B model, finetuned on the Open Assistant corpus, performs at a level comparable to ChatGPT according to the Vicuna benchmark. Since instruction-based finetuning is crucial for converting generic large language models into ChatGPT-like conversational agents, we anticipate that our method will facilitate widespread adoption of finetuning, particularly by researchers with limited computational resources. This advancement represents a significant step toward democratizing access to cutting-edge NLP tools. By bridging the gap between large technology firms and smaller groups operating with only consumer-grade hardware, \textsc{QLoRA} acts as a leveling force, expanding access and opportunity in the field.

An additional avenue for significant impact lies in enabling deployment on mobile devices. We argue that our \textsc{QLoRA} approach could serve as a key step toward making the finetuning of LLMs feasible on smartphones and other environments with limited computational resources. While it has been demonstrated previously that 7B parameter models can be executed on mobile phones, \textsc{QLoRA} represents the first methodology that supports the finetuning of these models directly on such devices. Our projections suggest that an iPhone 12 Plus is capable of finetuning on up to 3 million tokens overnight while being charged. Although finetuned 7B models do not attain the performance levels of ChatGPT, we consider their output sufficient to enable innovative use cases that were previously unviable due to privacy concerns or insufficient LLM performance. With \textsc{QLoRA}, it becomes possible to facilitate privacy-aware applications of LLMs, empowering users to maintain ownership and oversight over their data and models, while also reducing deployment barriers.

Nonetheless, it is important to acknowledge that finetuning presents risks as a dual-use technology, with potential for misuse and harm. The broader proliferation of LLMs entails well-documented challenges \citep{bommasani2021opportunities,bender2021dangers}, yet we contend that democratizing access to this rapidly advancing technology allows for more comprehensive and independent evaluation than would be possible if LLM capabilities were restricted to major corporations that do not disclose models or source code for external scrutiny.

In summary, we anticipate that \textsc{QLoRA} will substantially enhance the accessibility and convenience of finetuning high-quality LLMs, contributing positively by lowering barriers to their adoption.